\section{Conclusion}
\label{sec:conclusion}

This work studies how audited demonstrations can improve the next round of
real-robot teleoperation data collection. Successful nominal segments constrain
the model to preserve adequate human input, while correction intervals provide
local recovery behavior in difficult states. A visual-conditioned filter then
predicts a candidate expert action and a continuous assistance gain, composing
them through a bounded residual, gain-rate constraint, and safety projection.

The evaluation extends beyond offline action fitting to the collection process
itself. A frozen reference audit set, independent round-wise review, and a fixed
admission rule maintain comparable standards across iterations. Valid
demonstrations per operator minute measure collection efficiency, and downstream
imitation learning measures whether admitted episodes retain training value.
This framework makes the link between auditing, assistance, and recollection
explicit and experimentally testable for visually observable fine manipulation.
